\BourbakiExercisesSection

\begin{enumerate}
\def\labelenumi{\arabic{enumi}.}
\item
  Show that if E is a vector space \(\neq\{0\}\) over a field K with an infinity of elements, the set of generating systems of E is not inductive for the order relation \(\supset\) (form a decreasing sequence \((S_{n})\) of generating systems of E such that the intersection of the \(S_{n}\) is empty).
\item
  Let K be a field, L a subfield of K, I a set and \((x_{i})_{1\leqslant i\leqslant m}\) a free family of vectors in \(K_{s}^{I}\) such that the coordinates of each of the \(x_{i}\) belong to L. Let V be the vector subspace of \(K_{s}^{I}\) generated by the \(x_{i}\) ; show that there exists a finite subset J of I, with m elements, such that for every vector \(z=(\zeta_{\alpha})_{\alpha\in\mathbb{I}}\) of V, all the \(\zeta_{\alpha}\) belong to the subfield of K generated by L and the m elements \(\zeta_{\beta}\) where \(\beta\in J\) . (By applying Corollary 3 of no. 5, to the space \(L_{s}^{I}\) , show that it can be reduced to the case where there exist m indices \(\beta_{i}\in I\) ( \(1\leqslant i\leqslant m\) ) such that \(\Pr_{\beta_{i}}(x_{j})=\delta_{ij}\) ).
\item
  \begin{enumerate}
  \def\labelenumii{(\alph{enumii})}
  \tightlist
  \item
    Let G be a group and M an infinite subset of G. Show that if \(G'\) is the subgroup of G generated by M, then \(\text{Card}(G') = \text{Card}(M)\) .
  \end{enumerate}
\end{enumerate}

\begin{enumerate}
\def\labelenumi{(\alph{enumi})}
\setcounter{enumi}{1}
\item
  Let \(\mathbf{K}\) be a field and \(\mathbf{M}\) an infinite subset of \(\mathbf{K}\) . Show that if \(\mathbf{K}'\) is the subfield of \(\mathbf{K}\) generated by \(\mathbf{M}\) , then \(\operatorname{Card}(\mathbf{K}') = \operatorname{Card}(\mathbf{M})\) . (Consider two sequences \((\mathbf{A}_n)_{n \geqslant 0}, (\mathbf{P}_n)_{n \geqslant 0}\) of subsets of \(\mathbf{K}\) such that \(\mathbf{A}_0 = \mathbf{M}\) , \(\mathbf{P}_n\) is the multiplicative subgroup of \(\mathbf{K}^*\) generated by \(\mathbf{A}_n \cap \mathbf{K}^*\) and \(\mathbf{A}_{n+1}\) the additive subgroup of \(\mathbf{K}\) generated by \(\mathbf{P}_n\) ; then apply (a).)
\item
  Let K be a field and I an infinite set; show that \(\operatorname{Card}(\mathbf{K}) \leqslant \dim(\mathbf{K}_{s}^{\mathbf{I}})\) . (Reduce it to the case I = N and argue by reductio ad absurdum. Let B be a basis of \(K_{s}^{N}\) such that \(\operatorname{Card}(\mathbf{B}) < \operatorname{Card}(\mathbf{K})\) and let L be the subfield of K generated by the coordinates of all the elements of B; then \(\operatorname{Card}(\mathbf{L}) < \operatorname{Card}(\mathbf{K})\) . Form a sequence \((\xi_{n})_{n \geqslant 1}\) of elements of K such that for all n, \(\xi_{n}\) does not belong to the subfield of K generated by L and the \(\xi_{i}\) of index i \textless{} n; then apply Exercise 2 to obtain a contradiction, by considering the point \(x = (\xi_{n})\) of \(K_{s}^{N}\) .
\item
  Deduce from (c) that for every field K and every infinite set I,
\end{enumerate}

\[
\dim (\mathrm{K} _ {s} ^ {\mathrm{I}}) = (\operatorname{Card} (\mathrm{K})) ^ {\operatorname{Card} (\mathrm{I})}
\]

(``Erdös-Kaplansky theorem'').

4 Give an example of an infinite sequence \((\mathbf{F}_n)\) of vector subspaces of a vector space \(\mathbf{E}\) such that \(\operatorname{codim}\left(\bigcap_{n} \mathbf{F}_n\right) > \sum_{n} \operatorname{codim}(\mathbf{F}_n)\) . (Take \(\operatorname{codim}(\mathbf{F}_n) = 1\) for all \(n\) and use Exercise 3 \((d)\) .)

\begin{enumerate}
\def\labelenumi{\arabic{enumi}.}
\setcounter{enumi}{4}
\tightlist
\item
  Let \((\mathbf{H}_{\lambda})_{\lambda \in L}\) be a family of hyperplanes (passing through 0) of a vector space \(\mathbf{E}\) over a field \(\mathbf{K}\) , which is a covering of \(\mathbf{E}\) .
\end{enumerate}

\begin{enumerate}
\def\labelenumi{(\alph{enumi})}
\item
  Show that, if K is finite, then \(\operatorname{Card}(\mathbf{L}) \geqslant 1 + \operatorname{Card}(\mathbf{K})\) and, if K is infinite, \(\operatorname{Card}(\mathbf{L}) \geqslant \aleph_{0} = \operatorname{Card}(\mathbf{N})\) . (Prove by induction on r that if K has at least r elements, then E cannot be the union of r hyperplanes.) Show by examples that these inequalities cannot be improved without an additional hypothesis (when K is infinite, consider the space \(\mathbf{E} = \mathbf{K}_{s}^{(\mathbf{N})}\) ).
\item
  If E is finite-dimensional, show that \(\operatorname{Card}(\mathbf{L}) \geqslant 1 + \operatorname{Card}(\mathbf{K})\) . (Argue by induction on \(\dim(\mathbf{E})\) .) which is no longer true in infinite dim, consider \(K^{(N)}\)
\end{enumerate}

¶ 6. Let S be a non-empty finite subset of a vector K-space E, not containing 0. Suppose that there exists an integer \(k \geqslant 1\) with the following property: (*) for all \(x \in S\) , there is a partition of \(S = \{x\}\) into \(h\) free subsets with \(h \leqslant k\) .

For such a partition \((\mathbf{N}_i)_{1\leqslant i\leqslant h}\) and every sequence \((r_j)_{1\leqslant j\leqslant m}\) of integers of \(\{1,h\}\) a subspace \(\mathrm{F}_{r_1r_2\dots r_m}\) of \(\mathbf{E}\) is defined by induction on \(m\) as follows: \(\mathrm{F}_{r_1}\) is the subspace generated by \(\mathrm{N}_{r_1},\mathrm{F}_{r_1\dots r_p}\) the subspace generated by the intersection of \(\mathrm{F}_{r_1\dots r_{p-1}}\) and \(\mathrm{N}_{r_p}\) . Suppose that:

(**) for all \(x \in S\) and every partition \((N_i)_{1 \leqslant i \leqslant h}\) of \(S = \{x\}\) into at most \(k\) free subsets, there exists at least one sequence \((r_j)_{1 \leqslant j \leqslant m}\) such that

\[
x \notin \mathrm{F} _ {r _ {1} r _ {2} \dots r _ {m}}.
\]

\begin{enumerate}
\def\labelenumi{(\alph{enumi})}
\item
  Let \(n\) be the least of the integers \(m \geqslant 1\) for all possible choices of \(x \in S\) and let \((\mathbf{N}_i)_{1 \leqslant i \leqslant h}\) and \((r_j)_{1 \leqslant j \leqslant m}\) satisfy conditions (*) and (**). Show that of necessity \(n = 1\) . (Argue by reductio ad absurdum by assuming \(x \notin F_{r_1 r_2 \ldots r_n}\) and \(n > 1\) . Then necessarily \(x \in F_{r_n}\) . For simplicity write \(E_p = F_{r_1 r_2 \ldots r_p}\) for \(1 \leqslant p \leqslant n\) and \(x = \sum_i \alpha_i y_i\) with \(\alpha_i \in K\) and \(y_i \in N_{r_n}\) ; there is at least one \(y_i\) not belonging to \(E_{n-1}\) , denote it by \(y\) ; then write \(N = N_{r_n} - \{y\}\) , \(N_i' = N_i\) for \(1 \leqslant i \leqslant h\) and \(i \neq r_n\) , \(N_{r_n}' = N \cup \{x\}\) , so that \((\mathbf{N}_i')_{1 \leqslant i \leqslant h}\) is a partition of \(S - \{y\}\) consisting of free subsets. Define \(E_0' = E\) and for \(1 \leqslant p \leqslant n\) , \(E_p'\) as the subspace generated by the intersection of \(E_{p-1}'\) and \(N_{r_p}'\) . Now \(y \notin E_{n-1}\) ; let \(q\) be the least integer such that \(y \notin E_q\) ; show that \(E_q' \not\subset E_q\) . Let \(s\) be the least integer such that \(E_s' \not\subset E_s\) ; show that \(r_s = r_n\) ; deduce finally a contradiction from the relations \(y \in E_s, E_{s-1}' \subset E_{s-1}\) and \(E_s' \not\subset E_s\) .
\item
  Conclude from (a) that with the given hypotheses there exists a partition of S into h free subsets for an integer \(h \leqslant k\) .
\end{enumerate}

\begin{enumerate}
\def\labelenumi{\arabic{enumi}.}
\setcounter{enumi}{6}
\tightlist
\item
  Let S be a non-empty finite subset of a vector K-space E, not containing 0. Suppose that there exists an integer \(k \geqslant 1\) such that for every subset
\end{enumerate}

T of S, \(\text{Card}(T) \leqslant k.\text{rg}(T)\) . Show that there exists a partition of S into h free subsets for some integer \(h \leqslant k\) . (Prove that S satisfies condition ( \(^{**}\) ) of Exercise 6, arguing by induction on \(\text{Card}(S)\) and taking amongst all the subspaces \(F_{r_{1}\ldots r_{m}}\) one of those with the smallest possible dimension; if G is such a subspace and j an index for which \(\text{Card}(G \cap N_{j})\) is the least possible, prove that

\[
\begin{array}{r l} k. \operatorname{Card} (\mathrm{G} \cap \mathrm{N} _ {j}) & \leqslant \operatorname{Card} (\mathrm{G} \cap (\mathrm{S} - \{x \})) \\ & \leqslant \operatorname{Card} (\mathrm{G} \cap \mathrm{S}) \leqslant k. \operatorname{Card} (\mathrm{G} \cap \mathrm{N} _ {j}).) \end{array}
\]

\begin{enumerate}
\def\labelenumi{\arabic{enumi}.}
\setcounter{enumi}{7}
\item
  Let E be a vector space. Show that every endomorphism of E which is not a right divisor of 0 in End(E) is surjective (§ 2, Exercise 7).
\item
  Let K be a field.
\end{enumerate}

\begin{enumerate}
\def\labelenumi{(\alph{enumi})}
\item
  In the vector space \(\mathbf{E} = \mathrm{K}_s^{(\mathbf{Z})}\) , let \((e_n)_{n\in \mathbf{Z}}\) denote the canonical basis; let \(v\) be the automorphism of \(\mathbf{E}\) defined by the relations \(v(e_n) = e_{n + 1}\) for all \(n\in \mathbf{Z}\) . If \(u = 1_{\mathbb{E}} - v\) , show that \(u\) is an injective endomorphism of \(\mathbf{E}\) , but that \(\dim (\operatorname {Coker}(u)) = 1\) .
\item
  In the vector space \(\mathbf{E} = \mathrm{K}_s^{(\mathbb{N})}\) , let \((e_n)_{n\in \mathbb{N}}\) denote the canonical basis; let \(v\) be the endomorphism of \(\mathbf{E}\) defined by \(v(e_0) = e_0\) , \(v(e_n) = e_{n - 1} + e_n\) for \(n\geqslant 1\) . Show that \(v\) is an automorphism of \(\mathbf{E}\) and that, if \(u = 1_{\mathbb{E}} - v\) , \(u\) is a surjective endomorphism of \(\mathbf{E}\) such that \(\dim (\mathrm{Ker}(u)) = 1\) .
\item
  Let I be a non-empty set. Show that every automorphism u of the vector space \(E = K_{s}^{(I)}\) can be written as a difference u = v - w of two automorphisms of E, except when K has only two elements and I consists of a single element. (Note that it suffices to prove that one particular automorphism u is of the desired form; when K is a field with two elements, consider separately the cases where \(\text{Card}(I) = 2\) and \(\text{Card}(I) = 3\) .)
\end{enumerate}

¶ 10. Let E be a vector space over a field K. Show that every endomorphism u of E is an automorphism or can be written as u = v - w, where v and w are automorphisms of E. (Observe that u can always be replaced by \(s_{1} \circ u \circ s_{2}\) , where \(s_{1}\) and \(s_{2}\) are two automorphisms of E. Distinguish two cases, according to whether \(\dim(\text{Ker}(u)) \leqslant \dim(\text{Coker}(u))\) or \(\dim(\text{Ker}(u)) > \dim(\text{Coker}(u))\) ; when \(\text{rg}(u)\) is finite, then the first case always holds. When

\[
\dim (\operatorname{Ker} (u)) \leqslant \dim (\operatorname{Coker} (u)),
\]

it can be reduced to the case where \(\operatorname{Ker}(u) \cap \operatorname{Im}(u) = \{0\}\) ; if

\[
\dim (\operatorname{Ker} (u)) = \dim (\operatorname{Coker} (u)),
\]

it can also be assumed that \(\operatorname{Im}(u) + \operatorname{Ker}(u) = \mathbf{E}\) and then apply Exercise 9 (c). If \(\dim (\operatorname {Ker}(u)) < \dim (\operatorname {Coker}(u))\) , there is a supplementary subspace in \(\mathbf{E}\) of \(\operatorname {Ker}(u)\) of the form \(\mathbf{W} \oplus \operatorname{Im}(u)\) with \(\dim (\mathbf{W}) = \dim (\operatorname {Im}(u)) = \dim (\mathbf{E})\) ; apply the remark at the beginning with \(s_1(u(\operatorname {Im}(u))) = \operatorname{Im}(u)\) and use Exercise 9 (a) and (c). When \(\dim (\operatorname {Ker}(u)) > \dim (\operatorname {Coker}(u))\) , it can be reduced to the case where \(\operatorname{Ker}(u) \subset \operatorname{Im}(u)\) ; this time take \(s_1(\operatorname{Ker}(u)) = u(\operatorname{Ker}(u))\) and use the results of the preceding cases and Exercise 9.)

\begin{enumerate}
\def\labelenumi{\arabic{enumi}.}
\setcounter{enumi}{10}
\tightlist
\item
  Let E, F be two vector spaces over a field K and \(u: E \rightarrow F\) a linear mapping. Show that for every vector subspace V of F,
\end{enumerate}

\[
\dim (u ^ {- 1} (\mathbf {V})) = \dim (\mathbf {V} \cap \operatorname{Im} (u)) + \dim (\operatorname{Ker} (u)).
\]

\begin{enumerate}
\def\labelenumi{\arabic{enumi}.}
\setcounter{enumi}{11}
\tightlist
\item
  Let E, F be two vector spaces and u, v two linear mappings of E into F.
\end{enumerate}

\begin{enumerate}
\def\labelenumi{(\alph{enumi})}
\tightlist
\item
  Show that the following inequalities hold:
\end{enumerate}

\[
\begin{array}{c} \operatorname{rg} (v) \leqslant \operatorname{rg} (u + v) + \operatorname{rg} (u) \\ \operatorname{rg} (u + v) \leqslant \inf (\dim (\mathrm{E}), \dim (\mathrm{F}), \operatorname{rg} (u) + \operatorname{rg} (v)). \end{array}
\]

If E and F are finite-dimensional, show that \(\operatorname{rg}(u + v)\) can take any integer value satisfying the inequalities

\[
\left| \operatorname{rg} (u) - \operatorname{rg} (v) \right| \leqslant \operatorname{rg} (u + v) \leqslant \inf (\dim (\mathrm{E}), \dim (\mathrm{F}), \operatorname{rg} (u) + \operatorname{rg} (v)).
\]

\begin{enumerate}
\def\labelenumi{(\alph{enumi})}
\setcounter{enumi}{1}
\tightlist
\item
  Show that
\end{enumerate}

\[
\dim (\operatorname{Ker} (u + v)) \leqslant \dim (\operatorname{Ker} (u) \cap \operatorname{Ker} (v)) + \dim (\operatorname{Im} (u) \cap \operatorname{Im} (v)).
\]

\begin{enumerate}
\def\labelenumi{(\alph{enumi})}
\setcounter{enumi}{2}
\tightlist
\item
  Show that
\end{enumerate}

\[
\dim (\operatorname{Coker} (u + v)) \leqslant \dim (\operatorname{Coker} (u)) + \operatorname{rg} (v).
\]

(If \(\mathbf{V}\) is a supplementary subspace of \(\operatorname{Ker}(v)\) in \(\mathbf{E}\) , note that

\[
u (\operatorname{Ker} (v)) \subset \operatorname{Im} (u + v) \quad \text { and } \quad \dim (u (\mathbf {V})) \leqslant \operatorname{rg} (v).)
\]

\begin{enumerate}
\def\labelenumi{\arabic{enumi}.}
\setcounter{enumi}{12}
\tightlist
\item
  Let E, F, G be three vector spaces and \(u: \mathbf{E} \to \mathbf{F}\) , \(v: \mathbf{F} \to \mathbf{G}\) two linear mappings.
\end{enumerate}

\begin{enumerate}
\def\labelenumi{(\alph{enumi})}
\tightlist
\item
  Show that there exists a decomposition of E as a direct sum of \(\operatorname{Ker}(u)\) and two subspaces M, N such that \(\operatorname{Ker}(v \circ u) = \mathbf{M} \oplus \operatorname{Ker}(u)\) and
\end{enumerate}

\[
\operatorname{Im} (v \circ u) = v (u (\mathbf {N})).
\]

\begin{enumerate}
\def\labelenumi{(\alph{enumi})}
\setcounter{enumi}{1}
\item
  \(\dim(\operatorname{Ker}(u)) \leqslant \dim(\operatorname{Ker}(v \circ u)) \leqslant \dim(\operatorname{Ker}(u)) + \dim(\operatorname{Ker}(v))\) ; if \(\operatorname{Ker}(u)\) and \(\operatorname{Ker}(v)\) are finite-dimensional, show that \(\dim(\operatorname{Ker}(v \circ u))\) can take every integer value satisfying the above inequalities.
\item
  The following equalities hold:
\end{enumerate}

\[
\begin{array}{l} \operatorname{rg} (u) = \operatorname{rg} (v \circ u) + \dim (\operatorname{Im} (u) \cap \operatorname{Ker} (v)) \\ \operatorname{rg} (v) = \operatorname{rg} (v \circ u) + \operatorname{codim} _ {\mathbb {F}} (\operatorname{Im} (u) + \operatorname{Ker} (v)). \end{array}
\]

If F is finite-dimensional, show that \(\operatorname{rg}(v \circ u)\) can take any integer value satisfying the inequality

\[
\sup (0, \operatorname{rg} (u) + \operatorname{rg} (v) - n) \leqslant \operatorname{rg} (v \circ u) \leqslant \inf (\operatorname{rg} (u), \operatorname{rg} (v)).
\]

\begin{enumerate}
\def\labelenumi{\arabic{enumi}.}
\setcounter{enumi}{13}
\item
  Let E, F, G be three vector spaces and \(u: \mathbf{E} \to \mathbf{F}\) , \(w: \mathbf{E} \to \mathbf{G}\) two linear mappings. Show that if \(u(0) \subset w(0)\) , there exists a linear mapping \(v\) of F into G such that \(w = v \circ u\) (cf.~§ 2, Exercise 8).
\item
  Let E, F be two vector spaces over a field K and \(u: E \rightarrow F\) a linear mapping.
\end{enumerate}

\begin{enumerate}
\def\labelenumi{(\alph{enumi})}
\item
  The following conditions are equivalent: (1) \(\operatorname{Ker}(u)\) is finite-dimensional; (2) there exists a linear mapping \(v\) of \(\mathbf{F}\) into \(\mathbf{E}\) such that \(v \circ u = 1_{\mathbf{E}} + w\) , where \(w\) is an endomorphism of \(\mathbf{E}\) of finite rank; (3) for every vector space \(\mathbf{G}\) over \(\mathbf{K}\) and every linear mapping \(f: \mathbf{G} \to \mathbf{E}\) , the relation \(\operatorname{rg}(u \circ f) < +\infty\) implies \(\operatorname{rg}(f) < +\infty\) .
\item
  The following conditions are equivalent: (1) Coker(u) is finite-dimensional; (2) there exists a linear mapping v of F into E such that \(u \circ v = 1_{F} + w\) , where w is an endomorphism of F of finite rank; (3) for every vector space G over K and every linear mapping \(g: F \to G\) , the relation \(\operatorname{rg}(g \circ u) < +\infty\) implies \(\operatorname{rg}(g) < +\infty\) .
\end{enumerate}

\begin{enumerate}
\def\labelenumi{\arabic{enumi}.}
\setcounter{enumi}{15}
\tightlist
\item
  Let E, F be two vector spaces over K and \(u: E \to F\) a linear mapping. u is said to be of finite index if \(\operatorname{Ker}(u)\) and \(\operatorname{Coker}(u)\) are finite-dimensional and the number \(d(u) = \dim(\operatorname{Ker}(u)) - \dim(\operatorname{Coker}(u))\) is then called the index of u.
\end{enumerate}

Let G be a third vector space over K and \(v: F \rightarrow G\) a linear mapping. Show that if two of the three linear mappings u, v, \(v \circ u\) are of finite index, so is the third and \(d(v \circ u) = d(u) + d(v)\) (use Exercise 13 (a)).

\begin{enumerate}
\def\labelenumi{\arabic{enumi}.}
\setcounter{enumi}{16}
\tightlist
\item
  Let E, F, G, H be four vector spaces over K and \(u: E \rightarrow F\) , \(v: F \rightarrow G\) , \(w: G \rightarrow H\) three linear mappings. Show that
\end{enumerate}

\[
\operatorname{rg} (v \circ u) + \operatorname{rg} (w \circ v) \leqslant \operatorname{rg} (v) + \operatorname{rg} (w \circ v \circ u).
\]

\begin{enumerate}
\def\labelenumi{\arabic{enumi}.}
\setcounter{enumi}{17}
\item
  Let E, F be two vector spaces over a field K and u, v two linear mappings of E into F. For there to exist an automorphism f of E and an automorphism g of F such that \(v = g \circ u \circ f\) , it is necessary and sufficient that \(\operatorname{rg}(u) = \operatorname{rg}(v)\) , \(\dim(\operatorname{Ker}(u)) = \dim(\operatorname{Ker}(v))\) and \(\dim(\operatorname{Coker}(u)) = \dim(\operatorname{Coker}(v))\) .
\item
  Let E be a vector space over a field K.
\end{enumerate}

\begin{enumerate}
\def\labelenumi{(\alph{enumi})}
\item
  Show that every mapping \(f\) of \(\mathbf{E}\) into \(\mathbf{E}\) which is permutable with every automorphism of \(\mathbf{E}\) is of the form \(x \mapsto \alpha x\) , where \(\alpha \in \mathbf{K}\) (show first that for all \(x \in \mathbf{E}\) there exists \(\rho(x) \in \mathbf{K}\) such that \(f(x) = \rho(x)x\) , using the fact that \(f\) commutes with every automorphism of \(\mathbf{E}\) leaving \(x\) invariant).
\item
  Let \(g\) be a mapping of \(\mathbf{E} \times \mathbf{E}\) into \(\mathbf{E}\) such that, for every automorphism \(u\) of \(\mathbf{E}\) , \(g(u(x), u(y)) = u(g(x, y))\) for all \(x, y\) in \(\mathbf{E}\) . Show that there exist two elements \(\alpha, \beta\) of \(\mathbf{K}\) such that \(g(x, y) = \alpha x + \beta y\) on the set of ordered pairs \((x, y)\) of linearly independent elements of \(\mathbf{E}\) ; moreover there exists a mapping \(\phi\) of \(\mathbf{K} \times \mathbf{K}\) into \(\mathbf{K}\) such that \(g(\lambda x, \mu x) = \phi(\lambda, \mu)x\) for all \(x \in \mathbf{E}\) (same method). If further \(g(u(x), u(y)) = u(g(x, y))\) for every endomorphism \(u\) of \(\mathbf{E}\) , then \(g(x, y) = \alpha x + \beta y\) for all \(x, y\) in \(\mathbf{E}\) . Generalize to mappings of \(\mathbf{E}^n\) into \(\mathbf{E}\) .
\end{enumerate}

\begin{enumerate}
\def\labelenumi{\arabic{enumi}.}
\setcounter{enumi}{19}
\tightlist
\item
  Let E be a vector space.
\end{enumerate}

\begin{enumerate}
\def\labelenumi{(\alph{enumi})}
\item
  Show that if M and N are two vector subspaces of E and M', N' the orthogonals of \(\mathbf{M}\) and \(\mathbf{N}\) respectively in \(\mathbf{E}^*\) , the orthogonal of \(\mathbf{M} \cap \mathbf{N}\) is \(\mathbf{M}' + \mathbf{N}'\) (cf.~§ 2, Exercise 9 (a)).
\item
  If E is infinite-dimensional, show that there exist hyperplanes \(H'\) of \(E^{*}\) such that the subspace of E orthogonal to \(H'\) reduces to 0 (consider a hyperplane containing the coordinate forms corresponding to a basis of E).
\item
  Show that if \(\mathbf{E}\) is infinite-dimensional there exists an infinite family \((\mathbf{V}_i)\) of subspaces of \(\mathbf{E}\) such that, if \(\mathbf{V}_i'\) is the subspace of \(\mathbf{E}^*\) orthogonal to \(\mathbf{V}_i\) , the subspace of \(\mathbf{E}^*\) orthogonal to \(\bigcap_{i} \mathbf{V}_i\) is distinct from \(\sum_{i} \mathbf{V}_i'\) .
\item
  Deduce from (b) that if E is infinite-dimensional, there exists a decomposition \(V' \oplus W'\) of \(E^{*}\) as a direct sum, \(W'\) being finite-dimensional, for which the sum \(V + W\) of subspaces V, W of E orthogonal respectively to \(V'\) and \(W'\) is distinct from E.
\item
  For a subspace of E to be finite-dimensional, it is necessary and sufficient that its orthogonal in \(E^{*}\) be of finite codimension.
\end{enumerate}

\begin{enumerate}
\def\labelenumi{\arabic{enumi}.}
\setcounter{enumi}{20}
\item
  With the notation of Theorem 8 of no. 6, let \(u\) denote the linear mapping \(x \mapsto (\langle x, x_i^* \rangle)\) of \(\mathbf{E}\) into \(\mathbf{K}_s^{\mathbf{I}}\) and \(y_0 = (\eta_i)\) . \(\mathbf{K}_d^{(\mathbf{I})}\) is identified with a subspace of the dual of \(\mathbf{K}_s^{\mathbf{I}}\) under the canonical mapping of \(\mathbf{K}_d^{(\mathbf{I})}\) into its bidual; then show that if \(\mathbf{N}'\) is the kernel of \(^t u\) , the condition of Theorem 3 (no. 5) expresses the fact that \(y_0\) is orthogonal to the intersection of \(\mathbf{N}'\) and \(\mathbf{K}_d^{(\mathbf{I})}\) . When \(\mathbf{I}\) is infinite and the system (21) (no. 5) is of finite rank, show that this intersection is distinct from \(\mathbf{N}'\) (note that \(\mathbf{N}'\) is then of finite codimension).
\item
  Let E and F be two vector spaces over a field K and u a linear mapping of E into F. If V is a vector subspace of E and \(V'\) the orthogonal of V in \(E^{*}\) , show that the dual of \(u(V)\) is isomorphic to \({}^{t}u(F^{*})/(V'\cap{}^{t}u(F^{*}))\) . If \(W'\) is a subspace of \(F^{*}\) such that \({}^{t}u(W')\) is finite-dimensional and W is the orthogonal of \(W'\) in F, \({}^{t}u(W')\) is isomorphic to the dual of the space \(u(E)/(W\cap u(E))\) .
\item
  Show that for a linear mapping \(u\) of a vector space \(E\) into a vector space \(F\) to be such that \(\mathrm{rg}(t u) = \mathrm{rg}(u)\) , it is necessary and sufficient that \(\mathrm{rg}(u)\) be finite (cf.~Exercise 3 (d)).
\item
  Let E be a vector space of finite dimension n \textgreater{} 1 over a commutative field K; show that, unless n = 2 and K is a field with two elements, there exists no isomorphism \(\phi\) of E onto \(E^{*}\) depending only on the vector space structure of E. (Note that if \(\phi\) is such an isomorphism, then necessarily \(\langle x, \phi(y) \rangle = \langle u(x), \phi(u(y)) \rangle\) for x, y in E and for every automorphism u of E (Set Theory, IV, § 1, no. 5).)
\item
  \begin{enumerate}
  \def\labelenumii{(\alph{enumii})}
  \tightlist
  \item
    Let K be a field and L, M two sets; there are canonical inclusions \((\mathbf{K}_{s}^{\mathbf{L}})^{(\mathbf{M})}\subset(\mathbf{K}_{s}^{(\mathbf{M})})^{\mathbf{L}}\subset\mathbf{K}_{s}^{\mathbf{L}\times\mathbf{M}}\) . Show that for two of these vector spaces to be equal, it is necessary and sufficient that one of the sets L, M be finite.
  \end{enumerate}
\end{enumerate}

\begin{enumerate}
\def\labelenumi{(\alph{enumi})}
\setcounter{enumi}{1}
\tightlist
\item
  Let \((\mathbf{E}_{\lambda})_{\lambda \in L}\) be a family of right vector spaces over \(K\) and \((\mathbf{F}_{\mu})_{\mu \in M}\) a family of left vector spaces over \(K\) . For the canonical mapping (26) (no. 7) to be bijective, it is necessary and sufficient that one of the vector spaces \(\prod_{\lambda \in L} \mathbf{E}_{\lambda}, \prod_{\mu \in M} F_{\mu}\) be finite-dimensional (use (a)).
\end{enumerate}

\begin{enumerate}
\def\labelenumi{\arabic{enumi}.}
\setcounter{enumi}{25}
\tightlist
\item
  \begin{enumerate}
  \def\labelenumii{(\alph{enumii})}
  \tightlist
  \item
    Let E, F be two left vector spaces over a field K; for the canonical homomorphism \(E^{*} \otimes_{K} F \to \operatorname{Hom}_{K}(E, F)\) to be bijective, it is necessary and sufficient that one of the spaces E, F be finite-dimensional.
  \end{enumerate}
\end{enumerate}

\begin{enumerate}
\def\labelenumi{(\alph{enumi})}
\setcounter{enumi}{1}
\tightlist
\item
  Let E be a right vector space and F a left vector space over a field K. For the canonical homomorphism \(E \otimes_{K} F \to \operatorname{Hom}_{K}(E^{*}, F)\) to be bijective, it is necessary and sufficient that E be finite-dimensional (use Exercise 3(d)).
\end{enumerate}

\begin{enumerate}
\def\labelenumi{\arabic{enumi}.}
\setcounter{enumi}{26}
\tightlist
\item
  \begin{enumerate}
  \def\labelenumii{(\alph{enumii})}
  \tightlist
  \item
    Under the hypotheses of Proposition 16 (i) (no. 7), for the canonical homomorphism (27) (no. 7) to be bijective, it is necessary and sufficient that E or F be finite-dimensional (note that the image under (27) (no. 7) of an element of \(\mathrm{Hom}_{\mathbb{K}}(\mathbf{E},\mathbf{G})\otimes_{\mathbb{L}}\mathbf{F}\) maps E to a subspace of \(\mathbf{G}\otimes_{\mathbb{L}}\mathbf{F}\) contained in a subspace of the form \(\mathbf{G}\otimes_{\mathbb{L}}\mathbf{F}'\) , where \(\mathbf{F}'\) is a finite-dimensional subspace of \(\mathbf{F}\) ).
  \end{enumerate}
\end{enumerate}

\begin{enumerate}
\def\labelenumi{(\alph{enumi})}
\setcounter{enumi}{1}
\tightlist
\item
  Under the hypotheses of Proposition 16 (ii) (no. 7), for the canonical homomorphism (28) (no. 7) to be bijective, it is necessary and sufficient that one of the ordered pairs \((\mathbf{E}_1, \mathbf{E}_2), (\mathbf{E}_1, \mathbf{F}_1), (\mathbf{E}_2, \mathbf{F}_2)\) consist of finite-dimensional vector spaces (use (a) and Exercise 7 of §4).
\end{enumerate}

\begin{enumerate}
\def\labelenumi{\arabic{enumi}.}
\setcounter{enumi}{27}
\item
  Let \(\rho: \mathbf{K} \to \mathbf{A}\) be an injective homomorphism of a field \(\mathbf{K}\) into a ring \(\mathbf{A}\) and let \(\mathbf{E}\) be a left vector space over \(\mathbf{K}\) . For the canonical mapping \((\mathbf{E}^*)_{(\mathbf{A})} \to (\mathbf{E}_{(\mathbf{A})})^*\) to be bijective, it is necessary and sufficient that \(\mathbf{E}\) be finite-dimensional or that \(\mathbf{A}\) , considered as a right vector \(\mathbf{K}\) -space, be finite-dimensional (use Exercise 25 ( \(b\) )). When one of these two cases holds, there exists a canonical bijection \((\mathbf{E}^{**})_{(\mathbf{A})} \to (\mathbf{E}_{(\mathbf{A})})^{**}\) .
\item
  Let A be an integral domain, K its field of fractions, E an A-module and T(E) its torsion module.
\end{enumerate}

\begin{enumerate}
\def\labelenumi{(\alph{enumi})}
\item
  Show that every linear form on \(\mathbf{E}\) is zero on \(\mathrm{T}(\mathrm{E})\) , so that the duals \(\mathbf{E}^*\) and \((\mathrm{E} / \mathrm{T}(\mathrm{E}))^*\) are canonically isomorphic.
\item
  Show that the canonical mapping \((\mathbf{E}^{*})_{(\mathbf{K})}\to (\mathbf{E}_{(\mathbf{K})})^{*}\) is injective and that it is bijective when \(\mathbf{E}\) is finitely generated.
\item
  Give an example of a torsion-free A-module E such that \(E^{*} = \{0\}\) and \(E_{(K)} \neq \{0\}\) .
\end{enumerate}

¶ 30. Let A be an integral domain and E, F two A-modules. Show that if x is a free element of E and y a free element of F, then \(x \otimes y \neq 0\) in \(E \otimes_{A} F\) . (Reduce it to the case where E and F are torsion-free, then to the case where E and F are finitely generated and use Exercise 29 (b) to show that there is an A-bilinear mapping f of \(E \times F\) into K such that \(f(x, y) \neq 0\) .)

*31. Let \(\mathbf{K}_0\) be a commutative field, A the polynomial ring \(\mathrm{K}_0[\mathrm{X},\mathrm{Y}]\) in two indeterminates over \(\mathbf{K}_0\) , which is an integral domain, and E the ideal \(\mathrm{AX} + \mathrm{AY}\) of A. Show that in the tensor product \(\mathbf{E} \otimes_{\mathbf{A}} \mathbf{E}\) the element \(\mathbf{X} \otimes \mathbf{Y} - \mathbf{Y} \otimes \mathbf{X}\) is \(\neq 0\) and that \(\mathrm{XY}(\mathbf{X} \otimes \mathbf{Y} - \mathbf{Y} \otimes \mathbf{X}) = 0\) (consider the A-bilinear mappings of \(\mathbf{E} \times \mathbf{E}\) into the quotient module \(\mathbf{A}/\mathbf{E}\) ).*

\begin{enumerate}
\def\labelenumi{\arabic{enumi}.}
\setcounter{enumi}{31}
\item
  Extend the results of no. 10 to the case where A is a non-commutative ring admitting a field of left fractions K (I, §9, Exercise 15). (Note that if \(\xi_{i}\) ( \(1 \leqslant i \leqslant n\) ) are elements of K, there exists an \(\alpha \neq 0\) in A such that all the \(\alpha\xi_{i}\) belong to A.)
\item
  Let A be an integral domain. An A-module E is called divisible if for all \(x \in E\) and all \(\alpha \neq 0\) in A, there exists \(y \in E\) such that \(\alpha y = x\) .
\end{enumerate}

\begin{enumerate}
\def\labelenumi{(\alph{enumi})}
\item
  Let E, F be two A-modules. Show that if E is divisible, so is \(E \otimes_{A} F\) . If E is divisible and F is torsion-free, \(\operatorname{Hom}_{\mathbf{A}}(\mathbf{E}, \mathbf{F})\) is torsion-free. If E is a torsion A-module and F is divisible, then \(E \otimes_{A} F = \{0\}\) . If E is a torsion A-module and F is torsion-free, then \(\operatorname{Hom}_{\mathbf{A}}(\mathbf{E}, \mathbf{F}) = \{0\}\) .
\item
  Show that every injective A-module (§ 2, Exercise 11) is divisible and conversely that every divisible torsion-free A-module is injective (use criterion (δ) of 2, Exercise 11).
\end{enumerate}

\begin{enumerate}
\def\labelenumi{\arabic{enumi}.}
\setcounter{enumi}{33}
\tightlist
\item
  Let A be an integral domain, P a projective A-module, \(P'\) a projective submodule of P and \(j: P' \to P\) the canonical injection. Show that for every torsion-free A-module E, the homomorphism \(j \otimes 1: P' \otimes_{A} E \to P \otimes_{A} E\) is injective (reduce it to the case where E is finitely generated and embed E in a free A-module).
\end{enumerate}

¶ 35. (a) Let K be a field and E a (K, K)-bimodule. Suppose that the dimensions of E, as a left and right vector space over K, are equal to the same finite number n.~Show that there exists a family \((e_{i})\) of n elements of E, which is a basis for each of the two vector space structures on E. (Note that if \((b_{j})_{1\leq j\leq m}\) is a family of m\textless n elements of E, which is free for each of the two vector space structures on E and V is the left vector subspace and W the right vector subspace generated by \((b_{j})\) , either \(V+W\neq E\) , or \(V\cap CW\) and \(W\cap CV\) are non-empty; in the latter case, if \(y\in V\cap CW\) , \(z\in W\cap CV\) , \(y+z\) forms with the \(b_{j}\) a family of \(m+1\) elements which is free for each of the two vector space structures on E.)

\begin{enumerate}
\def\labelenumi{(\alph{enumi})}
\setcounter{enumi}{1}
\item
  Let F be a subbimodule of E, whose two vector space structures over K have the same dimension p \textless{} n.~Show that there exists a family \((e_{i})_{1 \leqslant i \leqslant n}\) of elements of E which is a basis of E for each of its vector space structures and such that \((e_{i})_{1 \leqslant i \leqslant p}\) is a basis of F for each of its two vector space structures (same method).
\item
  Let \((b_{j})_{1\leqslant j\leqslant n - 1}\) be a family of \(n - 1\) elements of \(\mathbf{E}\) , which is free for each of the two vector space structures on \(\mathbf{E}\) . Let \(\mathbf{V}\) (resp. W) be the hyperplane generated by the \((b_{j})\) in \(\mathbf{E}\) considered as a left (resp. right) vector space.
\end{enumerate}

Show that if \(\mathbf{V} \subset \mathbf{W}\) (resp. \(\mathbf{W} \subset \mathbf{V}\) ), then necessarily \(\mathbf{V} = \mathbf{W}\) (note that if \(a \notin \mathbf{W}\) , the set of \(\lambda \in \mathbf{K}\) such that \(\lambda a \in \mathbf{W}\) is a left ideal).

*36. (a) Let \(\mathbf{K}_0\) be a commutative field and \(\mathbf{K} = \mathbf{K}_0(\mathbf{X})\) the field of rational functions in one indeterminate over \(\mathbf{K}_0\) (Chapter IV). A (K, K)-bimodule structure is defined on K as follows: the product \(t.u\) of an element \(u \in \mathbf{K}\) by a left operator \(t \in \mathbf{K}\) is the rational function \(t(\mathbf{X})u(\mathbf{X})\) ; the product \(u.t\) of \(u\) by a right operator \(t \in \mathbf{K}\) is the rational function \(u(\mathbf{X})t(\mathbf{X}^2)\) ; the left (resp. right) vector K-space structure thus defined on K is then 1-dimensional (resp. 2-dimensional). Deduce examples of (K, K)-bimodules such that dimensions of the two vector space structures of such a bimodule are arbitrary integers.

\begin{enumerate}
\def\labelenumi{(\alph{enumi})}
\setcounter{enumi}{1}
\tightlist
\item
  Deduce from (a) an example of a (K, K)-bimodule E whose two vector space structures have the same dimension and such that there exists a sub-bimodule F of E whose two vector space structures do not have the same dimension.*
\end{enumerate}

¶ 37. (a) Let E be a left vector space (finite-dimensional or otherwise) and \(A_{i}\) ( \(1 \leqslant i \leqslant n\) ) vector subspaces of E. Suppose that, for every sequence \((a_{i})_{1 \leqslant i \leqslant n}\) of points of E such that \(a_{i} \in A_{i}\) for all i, the vector subspace of E generated by the \(a_{i}\) is of dimension \(\leqslant m\) (m an integer \(\leqslant n\) ). Then prove that there exists a subspace W of E of dimension \(h \leqslant m\) , containing \(h + (n - m)\) of the subspaces \(A_{i}\) . (Argue by induction on n.~Prove first that (for fixed n) attention may be confined to the case where m \textless{} n and \(\dim(A_{i}) \leqslant m\) for all i; it may also be assumed that \(A_{i} \neq \{0\}\) for all i and that the \(A_{i}\) are not all

of dimension 1. Then argue (for fixed \(n\) ) by induction on \(d = \sum_{i=1}^{m} \dim(A_i)\) . If for example \(\dim(A_n) \geqslant 2\) , consider a 1-dimensional subspace \(B_n\) of \(A_n\) of dimension 1 and apply the induction hypothesis to \(A_1, \ldots, A_{n-1}, B_n\) ; conclude that there exists a number \(k\) such that \(1 \leqslant k \leqslant m\) and a \(k\) -dimensional subspace U of E containing \(k\) of the \(A_i\) , for example \(A_1, \ldots, A_k\) ; replacing, if need be, \(k\) by an integer \(k'\) such that \(1 \leqslant k' \leqslant k\) , show also that it can be assumed that there exist vectors \(b_i \in A_i\) for \(1 \leqslant i \leqslant k\) such that \(b_1, \ldots, b_k\) form a basis of U; then project onto a supplementary subspace of U in E and use the induction hypothesis.)

\begin{enumerate}
\def\labelenumi{(\alph{enumi})}
\setcounter{enumi}{1}
\tightlist
\item
  Let K be a finite field and let E be an N-dimensional vector space over K. Let \(U_{i}\) ( \(1 \leqslant i \leqslant n\) ) be vector subspaces of E such that, for every subset H of \([1, n]\) , the intersection of the \(U_{i}\) such that \(i \in H\) has dimension \(\leqslant N - \text{Card}(H)\) . Show that the union of the \(U_{i}\) cannot be equal to E (argue by duality using (a)).
\end{enumerate}

\begin{enumerate}
\def\labelenumi{\arabic{enumi}.}
\setcounter{enumi}{37}
\item
  Let K be a commutative field of characteristic 0 and E a finite-dimensional vector space over K. Let \(p_{1}, \ldots, p_{m}\) be projectors of E onto vector subspaces of E such that \(1_{E} = p_{1} + p_{2} + \cdots + p_{m}\) . Show that E is the direct sum of the \(p_i(\mathbf{E})\) and that the \(p_i\) are pairwise permutable (take the traces of the two sides).
\item
  Let K be a field, L a subfield of K and \((\mathrm{K}_{\alpha})_{\alpha\in I}\) a left directed family of subfields of K such that \(L=\bigcap_{\alpha\in I}K_{\alpha}\) . Let E be a vector space over K and \((a_{i})_{1\leqslant i\leqslant m}\) a finite family of elements of E; show that if the family \((a_{i})\) is free over L, there exists an index \(\alpha\) such that \((a_{i})\) is free over \(K_{\alpha}\) . (Argue by induction on m-r, where r is the rank of the family \((a_{i})\) over K.)
\end{enumerate}

